\documentclass[11pt,a4paper]{article}
\usepackage[utf8]{inputenc}
\usepackage{amsmath,amssymb,amsfonts}
\usepackage{geometry}
\usepackage{hyperref}
\geometry{margin=1in}

\title{\textbf{Relativistic Time Signal Math: Lorentz Composition, Dynamic Auto-Scaling Bounds, and Phase Orbit Space Mapping}}
\author{Time Dilation DAW Research Group --- Version 0.0.1}
\date{August 2026}

\begin{document}
\maketitle

\section{Overview}
In Time Dilation DAW, time signals ($\gamma(t)$ velocity fields and local coordinate time $t_{\text{local}}$) operate as native audio-rate control streams. Objects such as \texttt{time.math\textasciitilde}, \texttt{time.scope}, and \texttt{time.xy} process, combine, and visualize relativistic time dilation vectors. This paper formalizes the Lorentz composition law for time signals and the dynamic auto-scaling algorithms for 2D phase visualizers.

\section{Lorentz Time Signal Composition Law}
When two independent relativistic time dilation sources with factors $\gamma_1 \ge 0.0$ and $\gamma_2 \ge 0.0$ are combined via relativistic addition (\texttt{time.math\textasciitilde} mode 0), their effective normalized velocities $v_1$ and $v_2$ relative to stationary coordinate time are:
\begin{equation}
v_1 = \gamma_1 - 1.0, \quad v_2 = \gamma_2 - 1.0
\end{equation}

By Einstein's Special Theory of Relativity, velocity addition in a spacetime manifold with speed-of-light bound $c = 4.0$ follows:
\begin{equation}
v_{\text{comp}} = \frac{v_1 + v_2}{1.0 + \frac{v_1 \cdot v_2}{c^2}}
\end{equation}

The composite relativistic time dilation factor $\gamma_{\text{comp}}$ is derived as:
\begin{equation}
\gamma_{\text{comp}} = 1.0 + v_{\text{comp}} = 1.0 + \frac{(\gamma_1 - 1.0) + (\gamma_2 - 1.0)}{1.0 + \frac{(\gamma_1 - 1.0)(\gamma_2 - 1.0)}{c^2}}
\end{equation}

This formula guarantees:
\begin{enumerate}
    \item \textbf{Causality Bound}: $v_{\text{comp}} < c$ for all finite $v_1, v_2 < c$.
    \item \textbf{Identity Element}: Combining any time signal $\gamma_1$ with un-dilated time $\gamma_2 = 1.0$ yields $\gamma_{\text{comp}} = \gamma_1$.
    \item \textbf{Symmetry}: $\gamma_{\text{comp}}(\gamma_1, \gamma_2) = \gamma_{\text{comp}}(\gamma_2, \gamma_1)$.
\end{enumerate}

\section{Dynamic Auto-Scaling Bounds Algorithms}

\subsection{Continuous Telemetry Scope Auto-Scaling (\texttt{time.scope})}
For a time signal sequence $y[i]$ over a sliding window history $i \in [0, K-1]$, the peak amplitude $M$ is tracked:
\begin{equation}
M = \max_{0 \le i < K} |y[i]|
\end{equation}
To prevent waveform clipping or off-screen rendering when $\gamma$ scales high ($\gamma \gg 1.0$), the scope peak scale $S[n]$ is smoothed via exponential moving average:
\begin{equation}
S[n] = S[n-1] + \lambda \cdot \left(\max(1.0, M \cdot 1.15) - S[n-1]\right)
\end{equation}
where $\lambda = 0.1$ is the smoothing coefficient and $1.15$ provides a $15\%$ visual headroom margin.

\subsection{Dual-Time 2D Phase Lissajous Orbital Bounds (\texttt{time.xy})}
For dual time signals $(x[i], y[i])$, the phase orbit radius $r[i]$ is:
\begin{equation}
r[i] = \sqrt{x[i]^2 + y[i]^2}
\end{equation}
The maximum radial extent $R_{\text{max}}$ over the orbital window is:
\begin{equation}
R_{\text{max}} = \max_{0 \le i < K} r[i]
\end{equation}
The 2D visual radius $R_{\text{scale}}[n]$ auto-scales smoothly to enclose the phase orbit:
\begin{equation}
R_{\text{scale}}[n] = R_{\text{scale}}[n-1] + \lambda \cdot \left(\max(1.0, R_{\text{max}} \cdot 1.15) - R_{\text{scale}}[n-1]\right)
\end{equation}
Normalized coordinates plotted on the screen vector canvas are $(\hat{x}[i], \hat{y}[i]) = (x[i] / R_{\text{scale}}, y[i] / R_{\text{scale}})$, ensuring full orbital visibility at any dilation magnitude.

\end{document}
